\documentclass[12pt]{article}
\usepackage[T1]{fontenc}
\usepackage[utf8]{inputenc}
\usepackage{lmodern}
\usepackage{textcomp}
\usepackage{enumitem}
\usepackage{geometry}
\geometry{margin=1in}
\begin{document}
\begin{center}
Answer Key - Astronomy - Graduate Level Assessment 23 \
\small (Answers are ordered exactly as in the question paper.)
\end{center}

\section*{Multiple Choice}
\begin{enumerate}[label=\arabic*.]
\item \textbf{Answer:} D
\\ \textit{Options:} A) They are the shattered remains of collisions between planets.; B) They are chunks of rock or ice that condensed long after the planets and moons had formed.; C) They are chunks of rock or ice that were expelled from planets by volcanoes.; D) They are leftover planetesimals that never accreted into planets.
\\ \textit{Note:} The correct answer is accurate because the Solar Nebular theory posits that the solar system formed from a rotating disk of gas and dust, where planetesimals were the building blocks of planets, and those that did not combine into planets became asteroids and comets, representing remnants of the early solar system.
\item \textbf{Answer:} C
\\ \textit{Options:} A) 1.7 million light years; B) 2.1 million light years; C) 2.5 million light years; D) 3.2 million light years
\\ \textit{Note:} The Andromeda Galaxy is located approximately 2.5 million light years from Earth, making it the nearest spiral galaxy to our Milky Way and a key reference point in astronomical studies of distance in the universe.
\item \textbf{Answer:} D
\\ \textit{Options:} A) stability of particles within rings.; B) gravitational effects at ring edges as the moons pass by.; C) gaps between rings.; D) Moons contribute to all of the above.
\\ \textit{Note:} Moons influence various astronomical phenomena, such as tidal forces on their parent planets, stabilization of axial tilt, and the potential for impacting climate and habitability, thus contributing to all the aspects mentioned in the question.
\item \textbf{Answer:} D
\\ \textit{Options:} A) it is about 20000 times stronger than Earth's magnetic field; B) it traps charged particles from Io's volcanoes in a "plasma torus" around the planet; C) it causes spectacular auroral displays at Jupiter's north and south poles; D) it switches polarity every 11 years
\\ \textit{Note:} Jupiter's magnetic field does not switch polarity every 11 years; instead, it has a stable magnetic field that differs from the solar magnetic cycle, which is characterized by the Sun's magnetic field reversing approximately every 11 years.
\item \textbf{Answer:} B
\\ \textit{Options:} A) nuclear fusion in the core; B) by contracting changing gravitational potential energy into thermal energy; C) internal friction due to its high rotation rate; D) chemical processes
\\ \textit{Note:} Astronomers believe Jupiter generates internal heat primarily through the process of gravitational contraction, where the planet's immense mass causes it to slowly compress under its own gravity, converting gravitational potential energy into thermal energy.
\end{enumerate}

\section*{True / False}
\begin{enumerate}[label=\arabic*.]
\setcounter{enumi}{5}
\item \textbf{Answer:} True
\\ \textit{Options:} A) True; B) False
\\ \textit{Note:} The nebular theory suggests that the formation of planets is influenced by temperature gradients in the protoplanetary disk, leading to the formation of rocky terrestrial planets closer to the Sun and gas giant jovian planets further away, thus predicting an unequal distribution rather than an equal number of each type.
\item \textbf{Answer:} True
\\ \textit{Options:} A) True; B) False
\\ \textit{Note:} Mercury's insufficient volcanic activity, which is a result of its small size and rapid cooling history, limits the planet's ability to release gases that could form a permanent atmosphere, thus supporting the statement as true.
\item \textbf{Answer:} False
\\ \textit{Options:} A) True; B) False
\\ \textit{Note:} A solar day on a planet is the time it takes for the Sun to return to the same position in the sky, which differs from the sidereal day length due to the planet's rotation and orbital movement, so if a solar day on Planet X lasts approximately 100 Earth days, its sidereal day must be shorter, making the statement false.
\item \textbf{Answer:} False
\\ \textit{Options:} A) True; B) False
\\ \textit{Note:} Type-Ia supernovae primarily produce optical and infrared emissions rather than significant X-ray emissions, as they result from the thermonuclear explosion of a white dwarf in a binary system, which does not generate the high-energy phenomena typical of X-ray binaries or supernovae of other types.
\item \textbf{Answer:} False
\\ \textit{Options:} A) True; B) False
\\ \textit{Note:} The azimuth of the north celestial pole (NCP) does not indicate geographic latitude in the sky because it represents a compass direction rather than a measure of angular distance from the celestial equator, which is necessary for determining latitude.
\end{enumerate}

\section*{Long Form Response}
\begin{enumerate}[label=\arabic*.]
\setcounter{enumi}{10}
\item \textbf{Answer:} Asteroids are characterized by their significant motion relative to background stars, which distinguishes them in sky surveys. Unlike stars that typically appear stationary due to their vast distances, asteroids exhibit noticeable movement across the celestial sphere over short observational periods. This motion, often described as a parallax effect, is a direct result of their proximity within the solar system, enabling astronomers to track their paths and predict future positions. Consequently, observational techniques must account for this relative motion, utilizing time-series photometry and astrometric measurements to accurately identify and characterize asteroids. Such unique characteristics necessitate specialized methodologies in data analysis, ensuring that asteroids are distinguished from the static backdrop of stars.
\\ \textit{Note:} The answer correctly identifies that asteroids are distinguished from background stars by their significant motion across the sky due to their proximity, which allows astronomers to observe and predict their trajectories effectively.
\item \textbf{Answer:} Hot Jupiters are gas giant exoplanets that form beyond the frost line, where temperatures are low enough for volatile compounds to condense into solid ice, facilitating the accumulation of mass necessary for planet formation. Their inward migration is thought to result from interactions with the protoplanetary disk, including gravitational interactions with surrounding material and other protoplanets, which can lead to changes in their orbits. Additionally, mechanisms such as tidal interactions with the central star and dynamical scattering can further drive these planets towards the star, often resulting in highly eccentric orbits. The eventual stabilization of their orbits can occur through interactions with other planets or the disk, leading to the observed close proximity of these massive planets to their host stars. Understanding this process provides insight into the diversity of planetary systems and the complex dynamics at play during planet formation.
\\ \textit{Note:} correct because it accurately describes the formation of hot Jupiters beyond the frost line, where solid ice contributes to mass accumulation, and it identifies the mechanisms, including gravitational interactions and tidal forces, that facilitate their inward migration toward their host stars.
\end{enumerate}

\vfill
\noindent\textit{End of Answer Key}
\end{document}